\begin{lemma}[Minimal-orbit conductors and leading stationary classes]
\label{mod:parameters-minimalorbitstationary}
Let $K/F$ be a ramified cyclic extension of prime degree, let $t$ be its
lower break, and put $T=t+1$.  Supply the residue-degree, integral
uniformizer, and monogenicity hypotheses used by
module~\ref{mod:ramification-normcharacters}.  For an exact
quasi-character datum $\chi$, Lean defines
\[
 \operatorname{MinimalOrbit}_{K/F}(\chi)
 \quad\Longleftrightarrow\quad
 \forall\mu\in S(K/F),\ \forall q,\quad
 m_F(\mu\chi)=q\Longrightarrow m_F(\chi)\le q.
\]
The quantifier is over the complete norm-character group, including the
trivial character.

For arbitrary exact conductor data $\chi_0$, Lean uses the finite
norm-character group and the least-conductor twist construction of
module~\ref{mod:ramification-normcharacters} to choose
$\mu_0\in S(K/F)$ and package
\[
 \chi=\mu_0\chi_0
\]
with its actual least conductor.  It proves that $\chi$ is minimal against
every further twist in the whole orbit.  With the proved finite instance
installed, error-term twist invariance is used in exactly the direction
\[
 \operatorname{Err}_{K/F}(\mu_0\chi_0)
 =\operatorname{Err}_{K/F}(\chi_0).
\]
This reduction assumes no First Main identity.

Now suppose that $\chi$ is minimal and put $m=m_F(\chi)$.  For every
$\mu\in S(K/F)$ Lean constructs the actual least conductor datum of
$\mu\chi$ and proves the complete piecewise formula
\[
 m_F(\mu\chi)=
 \begin{cases}
 m,&\mu=1,\\
 \max\{m,T\},&\mu\ne1.
 \end{cases}
\]
Thus the trivial character remains present but is never assigned the
positive conductor of a nontrivial norm character.  For $\mu\ne1$ the proof
is separated into the manuscript's three cases:
\begin{itemize}
\item if $m<T$, then $m_F(\mu\chi)=T>m$;
\item if $m=T$, then $m_F(\mu\chi)=T$, since whole-orbit minimality excludes
  the only possible conductor drop;
\item if $T<m$, then $m_F(\mu\chi)=m$.
\end{itemize}
The same exact-conductor theorem is exported for every nonidentity power
$\tau^j$.  No local class-field-theory input is used beyond the proved norm
character API.

The exact pullback formulas are also exposed with precisely this
whole-orbit minimality hypothesis.  If $\chi_K=\chi\circ N_{K/F}$ and
$m\le T$, then
\[
 m_K(\chi_K)=m.
\]
If $T\le m$, then
\[
 m_K(\chi_K)
 =[K:F]m-([K:F]-1)T.
\]
For additive data $\psi_K=\psi_F\circ\operatorname{Tr}_{K/F}$ this gives
the exact denominator identity
\[
 m_K(\chi_K)+n_K(\psi_K)
 =[K:F]\bigl(m+n_F(\psi_F)\bigr).
\]

Lean re-exports the full finite ramified orbit and its complete cyclic
enumeration by $\operatorname{Multiplicative}(\mathbf Z/[K:F]\mathbf Z)$.
The identity index maps to the trivial character and is a member of the
finite orbit.  Nonidentity indices satisfy the nontrivial-twist conductor
formula, and full products may be reindexed by the cyclic enumeration
without deleting the identity factor.

Assume finally that $m=T>1$.  Let $\psi$ have exact additive conductor
$n$, choose $\Gamma\in F^\times$ and $M\in\mathbf Z$ with
\[
 \operatorname{ord}_F(\Gamma)=M+n,
\]
and let $r$ carry the stationary-depth certificate
$T\le2r$ and $r+1\le T$.  For every nontrivial $\mu$, Lean proves the
honest quotient-class statement
\[
 \operatorname{lead}\bigl(
   \operatorname{St}_{\Gamma,r}(\mu)
   +\operatorname{St}_{\Gamma,r}(\chi)\bigr)\ne0
 \quad\text{in}\quad
 \mathfrak p_F^{M-T}/\mathfrak p_F^{M-T+1}.
\]
For a nonidentity natural power $\tau^j$, the common-layer power calculus
and stationary-class uniqueness prove, as a quotient equality,
\[
 \operatorname{St}_{\Gamma,r}(\tau^j)
 =j\operatorname{St}_{\Gamma,r}(\tau),
\]
and hence
\[
 \operatorname{lead}\bigl(
   j\operatorname{St}_{\Gamma,r}(\tau)
   +\operatorname{St}_{\Gamma,r}(\chi)\bigr)\ne0.
\]
For arbitrary explicitly supplied representatives $a,b$ of the two
quotient classes, Lean also derives
\[
 a+b\notin\mathfrak p_F^{M-T+1}.
\]
No representative is asserted to be canonical.

Before minimality is imposed, two exact conductor-$T$ characters can only
have product conductor $T$ or an actual conductor $q'<T$.  If
$1<q'<T$, Lean first constructs
\[
 \operatorname{St}_{\Gamma,\lceil q'/2\rceil}(\mu\chi)
 \in
 \mathfrak p_F^{M-q'}/
 \mathfrak p_F^{M-\lceil q'/2\rceil}
\]
at the new conductor and new depth, and only then maps it into the old
conductor-$T$ quotient, where its restriction equals the old sum.  If
$q'\le1$, Lean proves that no stationary depth exists and constructs no
stationary object.  Consequently the critical stationary theorem has no
tame $T=1$ instance, while it applies uniformly to the wild odd-prime and
wild quadratic branches; those later branches introduce no different
minimal-orbit boundary.

\LeanFile{LanglandsFirstMainLemma/Parameters/MinimalOrbitStationary.lean}
\lean{LanglandsFirstMainLemma.IsMinimalNormCharacterOrbitRepresentative}
\lean{LanglandsFirstMainLemma.ramifiedNormCharacterOrbitTwistData}
\lean{LanglandsFirstMainLemma.ramifiedNormCharacterOrbitTwistData_character}
\lean{LanglandsFirstMainLemma.minimalOrbitRepresentative}
\lean{LanglandsFirstMainLemma.minimalOrbitRepresentative_character}
\lean{LanglandsFirstMainLemma.minimalOrbitRepresentative_isConductor}
\lean{LanglandsFirstMainLemma.minimalOrbitRepresentative_mem_orbit}
\lean{LanglandsFirstMainLemma.minimalOrbitRepresentative_isMinimal}
\lean{LanglandsFirstMainLemma.errorTerm_minimalOrbitRepresentative}
\lean{LanglandsFirstMainLemma.minimalOrbit_nontrivialTwist_isConductor}
\lean{LanglandsFirstMainLemma.ramifiedNormCharacterOrbitTwistData_conductor_eq}
\lean{LanglandsFirstMainLemma.minimalOrbitTwistConductor}
\lean{LanglandsFirstMainLemma.ramifiedNormCharacterOrbitTwistData_conductor_eq_ite}
\lean{LanglandsFirstMainLemma.minimalOrbit_twist_conductor_eq_of_lt}
\lean{LanglandsFirstMainLemma.minimalOrbit_twist_conductor_eq_of_critical}
\lean{LanglandsFirstMainLemma.minimalOrbit_twist_conductor_eq_of_gt}
\lean{LanglandsFirstMainLemma.minimalOrbit_nontrivialPower_isConductor}
\lean{LanglandsFirstMainLemma.minimalOrbit_compNorm_conductor_eq_of_leCritical}
\lean{LanglandsFirstMainLemma.minimalOrbit_compNorm_conductor_eq_atOrAboveCritical}
\lean{LanglandsFirstMainLemma.minimalOrbit_compNorm_add_compTrace_conductor_eq}
\lean{LanglandsFirstMainLemma.minimalOrbitNormCharacterFinset}
\lean{LanglandsFirstMainLemma.one_mem_minimalOrbitNormCharacterFinset}
\lean{LanglandsFirstMainLemma.minimalOrbitZMod_zero_eq_one}
\lean{LanglandsFirstMainLemma.minimalOrbitZMod_twist_conductor_eq}
\lean{LanglandsFirstMainLemma.minimalOrbit_fullProduct_eq_zmodProduct}
\lean{LanglandsFirstMainLemma.minimalOrbit_stationary}
\lean{LanglandsFirstMainLemma.normCharacterPower_stationaryClass_eq_nsmul}
\lean{LanglandsFirstMainLemma.minimalOrbit_power_stationary}
\lean{LanglandsFirstMainLemma.minimalOrbit_stationary_representatives}
\lean{LanglandsFirstMainLemma.criticalNormCharacterTwist_conductor_eq_or_drop}
\lean{LanglandsFirstMainLemma.criticalNormCharacterTwist_stationaryClass_afterDrop}
\lean{LanglandsFirstMainLemma.normCharacterTwist_noStationaryClass_of_endpoint}
\leanok
\SourceText{Lemma lem:minimal-orbit-stationary.}
\uses{mod:delta-errorterm,mod:lamprecht-stationaryclasscalculus,
mod:ramification-normcharacters,mod:ramification-pullbackconductors}
\FormalizationContract
Minimality quantifies over the complete norm-character orbit.  All stored
twist conductors are actual least conductors.  The trivial character remains
in every finite enumeration and is handled separately.  The critical
noncancellation theorem is a nonzero statement in the leading lattice
quotient; its representative corollary quantifies over supplied
representatives and chooses no canonical field element.  A dropped product
class is constructed at its actual conductor and half-depth before any
comparison with the old layer, while conductors zero and one have no
stationary object.  Only the proved direction of error-term twist invariance
is used, and no downstream high-, low-, phase-, or wild-case theorem is an
input.
\end{lemma}
